\section{Order \& Right versus The Cathedral}

It's interesting that the title of Evola's work\footnote{\url{http://www.jrbooksonline.com/pdf\_books/revolt.pdf}} is almost a prescriptive piece of advice, or a command -``Revolt!'', \& that he advocates here a proper understanding of what a revolt would look like if it were legitimate: this is how a \emph{kshatriya} would act to maintain purity as he overthrew what was false. Yet there are several stringent preconditions. First, any half measures are rejected: one cannot revolt using revolutionary principles of any kind, which are doomed. Thus, the ``conservative'' option is stillborn and simply not an option -- one cannot build on that which is already sliding into the abyss. Secondly, the revolt is first of all as witness -- that is, one is impassive and immovable, standing still as all else goes over the cliff.

Thirdly, any revolt would have to include self-knowledge, as acting outside one's natural caste would be bound to cause inevitable complications which would compromise the outcome -- a peasant trying to lead kings would not come to any good end. This explains much of the Church's problem in resisting modernism : even spiritual brahmins are not authorized to assume the emperor's post. Since the West is engaged in a kind of global revolutionary enterprise, we are looking for the emperor's post, rather than a king who can re-establish centrality under an already recognized holy empire. It also explains the nature of anti-American sentiment that is brewing abroad -- even other decaying civilizations sense that resistance to the global trends somehow involve answering and ending American pre-eminence, rather than merely establishing some kind of vacuum.

What Moldbug has called The Cathedral\footnote{\url{http://unqualified-reservations.blogspot.com/2009/01/gentle-introduction-to-unqualified.html}} can rightly be viewed as a Demetrian impulse\footnote{\url{https://gornahoor.wordpress.com/2015/06/07/the-demetrian-race/}}; the liberal order in fact represents precisely what one could term a gynocratic drive towards the Nanny Church-State.

\begin{quotex}
The Cathedral, with its informal union of church and state, is positioned perfectly. It has all the advantages of being a formal arm of government, and none of the disadvantages. Because it formulates public policy, it is best considered our ultimate governing organ, but it certainly bears no responsibility for the success or failure of said policy. Moreover, it gets to program the little worm that is inserted in everyone's head, beginning at the age of five and going all the way through grad school.
\end{quotex}


Instead of regal legitimacy, the moral-sentimental impulses of a substitute order (bereft of real authority) lead by degenerate brahmins and warriors (mlecchas) will impose an alternate scheme which does not correspond to Cosmic Law.

The choice by Moldbug of the term ``Cathedral'' to describe this entity is particularly felicitous: about the time of the rise of the Romanesque/Gothic, according to Evola, the Papacy was in the process of humiliating the Empire in the investiture controversy -- Gregory made Henry IV stand in the snow at Canossa\footnote{\url{https://en.wikipedia.org/wiki/Investiture\_Controversy}}. Henry, true to the Evolian formula \& claim, had made his case in no uncertain terms: His letter ends, ``Henry, \textbf{\emph{king not through usurpation but through the holy ordination of God}}, to Hildebrand, at present not pope but false monk... I, Henry, king by the grace of God, with all of my Bishops, say to you, come down, come down, and be damned throughout the ages.'' In other words, contrary to the Cathedral, either of the Middle Ages, or of the Puritan era, or of our own era of secular hyper-Calvinism, the ordination of the holy Emperor flows, truly, from God directly, with neither people nor pope intervening. Evola speaks of the warrior-priest almost as if it was a fourth caste. It was to this order or class that the Templars belonged, who tellingly owed allegiance, not to the Ekklesia but to a Temple.

The cooperation of the American Church in endorsing and legitimizing the split between secular/sacred has been crucial, as American Christianity is not merely predominantly exoteric, but also (now) maternal and this-worldly: all transcendent proclivities have been curtailed, castrated, or exorcised. Thus, what really matters today for most church goers is excising the very elements which remain and remind them of what has been abandoned or lost. The alternative to revolution (transcendent hierarchy) has been rendered literally inconceivable. As Phillip Rieff\footnote{\url{http://www.amazon.com/Fellow-Teachers-Philip-Rieff/dp/0571106781}} would put it, any truly therapeutic reaction would be rendered still-born by instant attention and dissection by the existing organs of The Cathedral, including the media, the churches, and the government agencies. In privileging the ``spiritual'' over the profane, the original intention of transcendence has been lost, \& what has followed has been either the marginalizing of the Church in the face of secular Revolution, or, the trivializing of the secular and the spiritual in a false dichotomy that emasculates the transcendence of either. Moldbug\footnote{\url{http://www.slate.com/articles/technology/bitwise/2015/06/curtis\_yarvin\_booted\_from\_strange\_loop\_it\_s\_a\_big\_big\_problem.2.html}} makes clear how the political process functions under the secular-spiritual Cathedral:

\begin{quotex}
Without separation of church and state, it is easy be for a democracy to indulge itself in arbitrarily irresponsible misgovernment, simply by telling its bishops to inform their congregations that black is white and white is black. Thus misdirected, they are easily persuaded to support counterproductive policies which they wrongly consider productive. Union of church and state can foster stable iatrogenic misgovernment as follows. First, the church fosters and maintains a popular misconception that the problem exists, and the solution solves it. Secondly, the state responds by extruding an arm, agency, or other pseudopod in order to apply the solution. Agency and church are thus cooperating in the creation of unproductive or counterproductive jobs, as ``doctors.'' (my note: this is Phillip Rieff's ``Therapeutic Culture'') Presumably they can find a way to split the take. The root problem with a state church in a democratic state is that, to believe in democracy, one must believe that the levers of power terminate with the voters. But if your democracy has an effective state church, the actual levers of power pass through the voters, and go back to the church. The church teaches the voters what to think; the voters tell the politicians what to do. Naturally, it is easy for the politicians to short-circuit this process and just listen to the bishops. Thus the government has a closed power loop. With the church at its apex, of course.
\end{quotex}


Those of this Church who come to their senses (in some respect) have a difficult time getting their footing under them. For example, Oriana\footnote{\url{https://en.wikipedia.org/wiki/Oriana\_Fallaci}} Fallaci, in her confrontation\footnote{\url{http://www.freerepublic.com/focus/news/669794/posts}} with radical Islam, appealed, not to the traditional order of the West, but to the liberated serfs:

\begin{quotex}
Well, in my view America frees the plebes. Everyone is a plebe there. White, black, yellow, brown, purple, stupid, intelligent, poor, rich. Actually the rich are the most plebeian of all. Most of the time they're such boors! Crude, ill-mannered. You can tell immediately that they've never read \emph{Galateo}, that they've never had anything to do with refinement and good taste and sophistication. In spite of the money they waste on clothes, for example, they're so inelegant as to make the Queen of England look chic by comparison. But they are freed, by God. And in this world there is nothing stronger or more powerful than freed plebes. You will always get your skull cracked when you go up against the Freed Plebe. And they all got their skulls cracked by America: English, Germans, Mexicans, Russians, Nazis, Fascists, Communists. Even the Vietnamese got theirs cracked in the end, when they had to come to terms after their victory so that now when a former president of the United States goes there to visit they're in seventh heaven. ``Bienvenu, Monsieur le President, bienvenu!'' The problem is that the Vietnamese don't pray to Allah. It's going to be much harder to deal with the sons of Allah. Much longer and much harder. Unless the rest of the Western world stops peeing its pants. And starts reasoning a little and gives them a hand.
\end{quotex}


And it is disconcerting, and confusing at first, to see that ``what is left'' in the West seems to have attached itself to the symbols and ideals of the American Revolution, with all of its concomitant ideology. This is what the Kali Yuga looks like, \& it is why Evola is so prescient and precise in dealing with the possibilities offered to us in this time.

Ms. Fallaci alludes to something here important to Evola's case:

\begin{quotex}
The Italians have become such little lords. They vacation in Seychelles, come to New York to buy sheets at Bloomingdale's. \textbf{They're ashamed to be laborers and farmers, and won't be associated with the proletariat.} But those of whom I speak (immigrants/Muslims), what kind of laborers are they? What work do they do? In what way do they satisfy the demand for manual labor that the Italian ex-proletariat no longer supplies?
\end{quotex}


She had courage, culture, and breeding, but these qualities are not enough to do more than to notice the external processes of decay and to fight rear-guard, temporizing actions. At least she was paying attention. In the global West, technology and wealth has caused the abandonment of the castes, with the result that hordes of immigrants are required as virtual slave laborers. \textbf{\emph{Additionally, liberals like the engendered chaos}}. Moldbug offers further confirmation that Liberalism thrives and wins because of the \emph{libido dominandi}\footnote{\url{http://www.amazon.com/Libido-Dominandi-Liberation-Political-Control/dp/1587314657}} which Dante also argued was filling up Hell:

\begin{quotex}
One of the key reasons that intellectuals are fascinated by disorder, in my opinion, is the fact that disorder is an extreme case of complexity. And as you make the structure of authority in an organization more complex, more informal, or both -- as you fragment it, eliminating hierarchical execution structures under which one individual decides and is responsible for the result, and replacing them with highly fragmented, highly consensual, and highly process-oriented structures in which ten, twenty or a hundred people can truthfully claim to have contributed to the outcome, you increase the amount of power, status, patronage, and employment produced. Of course, you also make the organization less efficient and effective, and you make working in it a lot less fun for everyone -- you have gone from startup to Dilbert. This is Brezhnevian sclerosis, the fatal disease of organizations in a highly regulated environment. All work is guided by some systematic process, in which each rule was contributed by someone whose importance was a function of how many rules he added. In the future, we will all work for the government. Individually, this is the last thing your average intellectual wants to do, but it is the direction in which his collective acts are pushing us. In short: intellectuals cluster to the left, generally adopting as a social norm the principle of pas d'ennemis a gauche, pas d'amis a droit, because like everyone else they are drawn to power. The left is chaos and anarchy, and the more anarchy you have, the more power there is to go around. The more orderly a system is, the fewer people get to issue orders. The same asymmetry is why corporations and the military, whose system of hierarchical executive authority is inherently orderly, cluster to the right.
\end{quotex}


In other words, it is nothing other than a desire to usurp the higher spiritual castes that produces \emph{the strangely synchronized \& seemingly invincible movement of Leftism}. The same desire on the economic level has created the political vacuum being filled by immigration. In all cases, \textbf{it is the abandonment of individual \emph{ius} and collective \emph{fides} (Evola) which opens the door of chaos}. No one knows who the Emperor is, or (worse) would recognize him if he were to appear. The Cathedral denies that there is such a thing: the only Emperor (for them) is the hyper-individual, whose creation is ongoing and perversely polymorphous (Herbert Marcuse). It is doubtful if The Cathedral will be able to totally control the forces which they are playing around with. In fact, they are comfortable with using the chaos against their enemies, deploying it, so to speak, tactically.

It is highly interesting that in fact we are witnessing the investiture folly worked out all over again, yet with even more dire consequences. For another good historical analysis of this trend (but one with ideological problems) one should consult Rosenstock-Huessy's \emph{Out of Revolution}, which connects the dots between the various revolutions. We are now in an accelerated phase of this process, \& understanding the ``cloak of maliciousness'' which Liberalism covers its bare parts with will allow discernment of the ways Res Publica is becoming Res Pubica, and what seems likely to happen next. Moldbug's analysis concretely demonstrates how a false elite can benefit from ``power-sharing'' and revolution at the expense of the common good, \& at the same time get long term results they say they eschew. But it is Evola who touches the point with a needle. The good news is, as twilight sets in, that many who are latently able have not yet begun to fight, \& that there is a deeper law at work, capable of being discerned and connected with at the metaphysical level.


\flrightit{Posted on 2015-07-18 by Logres}

\begin{center}* * *\end{center}

\begin{footnotesize}
\begin{sffamily}
\texttt{Mark Citadel on 2015-07-23 at 14:07 said:}

``Evola speaks of the warrior-priest almost as if it was a fourth caste. It was to this order or class that the Templars belonged, who tellingly owed allegiance, not to the Ekklesia but to a Temple.''

The Templars are very interesting. I have spoken in private communiques with Orthodox Christians in Russia and Belarus about the need for an organization that can operate outside the purview of the official church and powers that be, almost like a mimic of what Hezbollah did in Lebanon, but in the Orthodox world. This would of course be an entirely anti-Liberal venture and a preparatory force for the chaotic `twilight' we are entering into.

I really liked this essay. I am admittedly no scholar when it comes to Moldbug, but you compliment his work nicely here.

\hfill

\texttt{Logres on 2015-08-02 at 09:12 said:}

Moldbug is purely a comedic ironist, with little pretensions to anything beyond analyzing or ``auditing'' the USG (United States government). However, he does a good job of clinically taking apart the Cathedral \& exposing the pretensions of the man behind the curtain. I do think he is right about how power-sharing is driving the evolution of hyper-democracy and Revolution, and haven't seen anyone else pull apart the mechanics at that level. Of course, this is seen as ``good'' or even ultimate good, which is why it is invisible. I'm reading James Fitzjames Stephens' critique of JS Mill's On Liberty, \& he is saying a very similar thing about the secular government's potential to become effectively a Church; additionally, it clarifies that Mill lies at the root of a great deal of the modern Cathedral's religion. Mill was a powerful writer, and a poor thinker -- the distinction between ``other regarding acts'' and ``self regarding acts'' is like (as Stephens says) asking which acts of ours take place in space, and which in time: it's artificial, designed to produce a rhetorical conclusion. Stephens clarifies the more traditional consensus that held in the British Empire, in which certain forms of coercion were legitimate, and were connected to moral and philosophical principles.


\hfill

\end{sffamily}
\end{footnotesize}

